This paper investigates the coupling between geometric stress concentrations and nonlinear structural instabilities in
compressible hyperelasticity, focusing on a corner subject to mixed Dirichlet--Neumann boundary conditions. Existing
literature treats fracture, cavitation, and flat-surface creasing in isolation, leaving the mixed-boundary corner
unresolved. We establish a dichotomy in the asymptotic behaviour of the linearised tangent operator, dictated by the
loading direction and the volumetric energy penalty. In tension, a logarithmic volumetric penalty yields a benign
singularity where the tangent operator remains strongly elliptic and Fredholm, preserving the standard Kondratiev edge
exponents. Under compression, the local volume contraction drives an algebraic blow-up of the operator norm. We show
that this breakdown is a structural bifurcation rather than a material loss of ellipticity. Using a tangent-cone
blow-up reduction, we prove that a localised surface creasing instability preempts pointwise material collapse. An
imperfect bifurcation analysis demonstrates that the quadratic energy coefficient vanishes by symmetry, whereas the
subcritical cubic coefficient is tip-convergent yet explicitly dependent on the global stress-intensity factor, an
outcome of the characteristic edge exponent satisfying $\alpha > 1/2$. The theoretical framework is supported by a
scale-covariant Mellin--Chebyshev spectral formulation that resolves the singular exponents and buckling thresholds
with exponential accuracy.